% ---
% titre: Périmètre et aire d'une figure complexe
% theme: GM
% chapitre: Lignes et surfaces
% sous_chapitre: Lignes et surfaces
% niveau: [9e]
% type: EVAL
% tags: [c-perimetre, c-aire, c-figure-complexe, p-question-TS]
% difficulte: 2
% corrige: true
% ---

\question
Voici une figure $ABCDEFG$. Sachant que $ABCD$ est un parallélogramme: 

\vfill
\begin{center}
\begin{tikzpicture}[scale=1.5]
    \tkzDefPoints{0/0/A,8/0/D,2.3/4/B,10.3/4/C,3/0/G,7/0/E,4/3/F,4/0/H,-1/0/I, -1/4/J}
    \tkzDrawSegments(A,B B,C C,D D,E E,F F,G G,A)
    \tkzDrawSegments[dashed](E,G)
    %
    \tkzLabelPoints[below](A,D,E,G)
    \tkzLabelPoints[above](B,F,C)
    %
    \tkzDrawSegment[draw=none, dim={8\,m,20pt,midway}, dim style/.append style={black,dashed},dim fence style/.append style={draw=none}](B,C)
    \tkzDrawSegment[draw=none, dim={{4,6\,m},15pt,midway}, dim style/.append style={black,dashed},dim fence style/.append style={draw=none}](C,D)
    \tkzDrawSegment[draw=none, dim={{4,2\,m},-15pt,midway}, dim style/.append style={black,dashed},dim fence style/.append style={draw=none}](E,F)
    \tkzDrawSegment[draw=none, dim={{3,5\,m},-15pt,midway}, dim style/.append style={black,dashed},dim fence style/.append style={draw=none}](F,G)
    \tkzDrawSegment[draw=none, dim={3\,m,0pt,midway}, dim style/.append style={black,dashed},dim fence style/.append style={draw=none}](F,H)
    \tkzDrawSegment[draw=none, dim={4\,m,20pt,midway}, dim style/.append style={black,dashed},dim fence style/.append style={draw=none}](E,G)
    \tkzDrawSegment[draw=none, dim={4\,m,-20pt,midway}, dim style/.append style={black,dashed},dim fence style/.append style={draw=none}](I,J)
    %
    \tkzMarkRightAngle[size=0.3,draw,dashed](F,H,A)
\end{tikzpicture}
\end{center}

\begin{parts}
\vfill\part[3\half]
Calcule le périmètre de la figure ABCDEFG.

\vspace{10pt}
{\footnotesize\raggedleft Espace pour ta démarche et tes calculs\par}
\begin{solutionorgrid}[200pt]
Comme $ABCD$ est un parallélogramme: 

\vspace{5pt}
$CD \cong AB = 4,6\,\text{m}$ et $AD \cong BC = 8\,\text{m}$ \hfill \textbf{(1pt)}

\vspace{10pt}
$DE + GA = AD - GE = 8 - 4 = 4\,m$ \hfill \textbf{(1pt)}

\vspace{5pt}
$P_{ABCDEFG} = AB+BC+CD+(DE+GA)+EF+FG = 4,6+8+4,6+4+4,2+3,5 = 28,9\,m$ \hfill \textbf{(1pt)}
\end{solutionorgrid}

\begin{tcolorbox}[enhanced jigsaw, interior hidden, colframe=box, parbox=false]%
Réponse: 
\sol{\vspace{20pt}}{Le périmètre de la figure vaut 28,9\,m. \textbf{(0,5pt)}}
\end{tcolorbox}

\newpage\part[3\half]
Calcule l'aire de la figure ABCDEFG.

\vspace{10pt}
{\footnotesize\raggedleft Espace pour ta démarche et tes calculs\par}
\begin{solutionorgrid}[200pt]
La figure $ABCDEFG$ correspond au parallélogramme $ABCD$ dont on soustrait le triangle $GFE$.

\vspace{5pt}
$A_{ABCD} = AD \cdot h = 8 \cdot 4 = 32\,m^2$ \hfill \textbf{(1pt)}

\vspace{5pt}
$A_{GFE} = \dfrac{GE \cdot h}{2} = \dfrac{4\cdot 3}{2} = 6\,m^2$ \hfill \textbf{(1pt)}

\vspace{5pt}
$A_{ABCDEFG} = 32 - 6 = 26\,m^2$ \hfill \textbf{(1pt)}
\end{solutionorgrid}

\begin{tcolorbox}[enhanced jigsaw, interior hidden, colframe=box, parbox=false]%
Réponse: 
\sol{\vspace{20pt}}{L'aire de la figure vaut 26\,m$^2$. \textbf{(0,5pt)}}
\end{tcolorbox}
\end{parts}